\chapter{Неограниченный родитель расширения Тёплица и гейт степени}

Bott-сравнение доказало вещественный класс пятнадцать, но использовало
ограниченный фредгольмов представитель. Теперь проверяется, выводится ли
пространство Харди и сдвиг из канонического неограниченного оператора, а
также достаточно ли этого для полного KO6-цикла.

\section{Стандартный неограниченный класс расширения}

Пусть
\begin{equation}
 \mathcal H=\ell^2(\mathbb Z)\otimes q_0\mathbb C^{105},
 \qquad
 Ne_n=ne_n,
 \qquad
 Ue_n=e_{n+1}.
 \label{eq:v5-unbounded-number-shift}
\end{equation}
Тогда на конечных линейных комбинациях базисных векторов
\begin{equation}
 [N,U]=U,
 \qquad
 [N,U^*]=-U^*.
 \label{eq:v5-unbounded-commutators}
\end{equation}
Резольвента $(1+N^2)^{-1/2}$ компактна, поскольку её собственные значения
$(1+n^2)^{-1/2}$ стремятся к нулю.

Бурн, Келлендонк и Ренни строят для вещественного или комплексного
расширения Тёплица неограниченный Kasparov-модуль
\begin{equation}
 \bigl(C(S^1)\widehat\otimes Cl_{0,1},
 \ell^2(\mathbb Z)\widehat\otimes\Lambda^*\mathbb R,
 N\otimes\gamma_{\rm ext}\bigr),
 \label{eq:v5-unbounded-literature-cycle}
\end{equation}
представляющий класс расширения; см. предложение 3.3 в
\path{arXiv:1604.02337v3}. Общий обзор расширений Тёплица и их
неограниченных представителей дан в \path{arXiv:1911.05823}.

\section{Поляризация Харди теперь выведена}

Спектральный проектор неотрицательной части $N$ равен
\begin{equation}
 P=\chi_{[0,\infty)}(N).
 \label{eq:v5-unbounded-spectral-projection}
\end{equation}
Его образ есть
\begin{equation}
 P\ell^2(\mathbb Z)=\ell^2(\mathbb N_0)\simeq H^2(S^1).
\end{equation}
Поэтому пространство Харди больше не является отдельно выбранным
аналитическим сектором: оно получается функциональным исчислением
канонического оператора числа.

Компрессии двух символов имеют вид
\begin{equation}
 PUP=S,
 \qquad
 PU^*P=S^*.
 \label{eq:v5-unbounded-compressions}
\end{equation}
На бесконечномерном образе $P$ выполнено
\begin{align}
 S^*S&=I,& SS^*&=I-|e_0\rangle\langle e_0|,\
 \operatorname{ind}S&=-1,& \operatorname{ind}S^*&=+1.
\end{align}
После тензорного умножения на $q_0$ ранга пятнадцать возвращаются
индексы $(-15,+15)$ и дефектный проектор
\begin{equation}
 |e_0\rangle\langle e_0|\otimes q_0
\end{equation}
ранга пятнадцать и веса $1/7$.

\begin{proposition}[Внутреннее происхождение поляризации]
Ограниченный коэффициентный цикл Тёплица является компрессией символов
$U,U^*$ спектральным проектором неограниченного оператора $N$. Поэтому
его пространство Харди, компактный дефект и индексы выводятся из одного
неограниченного аналитического класса расширения.
\end{proposition}

\section{Вещественный обмен}

Пусть $K$ --- комплексное сопряжение функций на окружности. В базисе
$e_n(z)=z^n$ оно действует как $Ke_n=e_{-n}$, поэтому
\begin{equation}
 KNK^{-1}=-N,
 \qquad
 KUK^{-1}=U^*.
 \label{eq:v5-unbounded-real-exchange}
\end{equation}
Вместе с уже доказанным KO6-каркасом двух сопряжённых копий это даёт
антилинейный обмен двух символов и двух ориентаций. Таким образом,
аналитическая конструкция совместима с проектной Real-инволюцией.

Следует, однако, различать обмен поляризаций и простое дополнение
проектора. При выбранном $P=\chi_{[0,\infty)}(N)$ нулевая мода принадлежит
положительной части, а $KPK^{-1}=\chi_{(-\infty,0]}(N)$. Пересечение двух
поляризаций равно линии $e_0$. Именно эта граничная линия создаёт
компактный проектор индекса; её нельзя удалить симметричным конечномерным
усечением.

\section{Степень расширения и соглашение о знаке}

Здесь возникает решающий контроль против ложного вывода. Модуль
\eqref{eq:v5-unbounded-literature-cycle} представляет класс
\begin{equation}
 [\mathrm{ext}]\in KKO_1(C(S^1),\mathbb R)
\end{equation}
и несёт действие $Cl_{0,1}$. Он является неограниченным родителем
расширения Тёплица, но сам по себе не является классом степени шесть.

Номер $KKO_1$ в обозначениях неограниченного цикла нельзя напрямую
складывать с индексом ковариантной таблицы $KO_n$, использованной в
Bott-сравнении. В стандартной точной последовательности вещественной
$K$-теории класс расширения задаёт карту
\begin{equation}
 \partial_n:KO_n(A)\longrightarrow KO_{n-1}(B).
\end{equation}
Поэтому для получения доказанного ранее выходного класса в $KO_6$
входной вещественный символ должен иметь степень семь:
\begin{equation}
 7-1=6\pmod 8.
 \label{eq:v5-unbounded-degree-ledger}
\end{equation}
Требуется построить
\begin{equation}
 [u_{\mathbb R},q_0]\in KO_7(C^*_{\mathbb R}(\mathbb Z)\otimes
 M_{105}(\mathbb C)_{\mathbb R})
\end{equation}
так, чтобы его Kasparov-произведение с $[\mathrm{ext}]$ имело
комплексификацию $(-15,+15)$.

\begin{theorem}[Разделение аналитического и KO6-родителя]
Оператор числа $N$ канонически выводит расширение Тёплица, поляризацию
Харди и коэффициентные индексы. Это закрывает аналитический
неограниченный разрыв. Полный неограниченный KO6-цикл ещё не закрыт,
поскольку степень семь обменного Real-символа пока не предъявлена.
\end{theorem}

Предыдущий результат $[T_{\mathbb R}]=15$ от этого не понижается: он уже
доказан инъективным Bott-сравнением ограниченного класса. Новый разрыв
касается его факторизации как конкретного неограниченного
Kasparov-произведения.

\begin{nogocriterion}[Запрет смешения степеней]
Нельзя объявлять модуль числа $N$ полным KO6-родителем только потому, что
его компрессия имеет нужные индексы. До построения входного Real-символа
степени семь доказан лишь аналитический класс расширения.
\end{nogocriterion}

\begin{protocol}
\begin{enumerate}
 \item неограниченный оператор числа $N$: \textbf{построен};
 \item ограниченность $[N,U]$: \textbf{выполнена};
 \item компактность резольвенты: \textbf{выполнена};
 \item проектор Харди как $\chi_{[0,\infty)}(N)$: \textbf{выведен};
 \item компрессии $PUP$ и $PU^*P$: \textbf{$S$ и $S^*$};
 \item коэффициентные индексы: \textbf{$-15/+15$};
 \item дефектный вес: \textbf{$1/7$};
 \item вещественный обмен $N\leftrightarrow-N$, $U\leftrightarrow U^*$:
 \textbf{выполнен};
 \item степень класса расширения: \textbf{один};
 \item входной Real-символ степени семь: \textbf{не построен};
 \item полный неограниченный KO6-родитель: \textbf{не закрыт};
 \item класс пятнадцать ограниченного цикла: \textbf{сохраняется};
 \item физическое действие: \textbf{не получено};
 \item следующий гейт: \textbf{обменный Real-символ степени семь}.
\end{enumerate}
\end{protocol}

Машинный сертификат сохранён в
\path{s2t_v5_real_toeplitz_unbounded_parent_cycle_gate_results.json}.